% Generated draft: requires mathematical and source review.
\documentclass[11pt]{article}
\usepackage[margin=25mm]{geometry}
\usepackage{fontspec}
\setmainfont{Latin Modern Roman}
\usepackage{amsmath,amssymb,mathtools,unicode-math}
\setmathfont{Latin Modern Math}
\usepackage{enumitem,needspace,xurl,hyperref}
\setlength{\parindent}{0pt}
\setlength{\parskip}{5pt}
\setlength{\emergencystretch}{3em}
\newcommand{\N}{\mathbb{N}}
\newcommand{\Z}{\mathbb{Z}}
\newcommand{\Q}{\mathbb{Q}}
\newcommand{\R}{\mathbb{R}}
\newcommand{\C}{\mathbb{C}}
\newcommand{\F}{\mathbb{F}}
\newcommand{\A}{\mathbb{A}}
\newcommand{\PP}{\mathbb{P}}
\newcommand{\E}{\mathbb{E}}
\newcommand{\eps}{\varepsilon}
\newcommand{\dd}{\,\mathrm{d}}
\begin{document}
\section*{UnsolvedMath ID: 30003289}
\textbf{Source problem code:} catalogue entry
\subsection{Definitions and mathematical statement}
An integral homology \(3\)-sphere is a closed connected \(3\)-manifold \(M\) with \(H_i(M;\Z)\cong H_i(S^3;\Z)\) for every \(i\geq0\). A genus-\(g\) Heegaard splitting expresses \(M\) as the union of two genus-\(g\) handlebodies along their common boundary. Stabilization of such a splitting adds the standard genus-one splitting of \(S^3\).

The supplied record refers to sets \(V_k(3)\) as Johnson-filtration images associated with stable Heegaard splittings of integral homology \(3\)-spheres. It does not define their domain, target, image map, or stabilization identifications. Retaining this notation, the first requested problem is to describe \(V_k(3)\) and determine whether
\[
\forall k\in\Z_{\geq1},\qquad V_k(3)=V_1(3).
\]

The second requested problem uses a group \(\pi\), subgroups \(\Gamma_{k+1}\), and groups denoted \(A_{g,1}(k)\), \(B_{g,1}\), and \(AB_{g,1}\), whose elements act on \(\pi\). For a normal subgroup \(\Gamma_{k+1}\trianglelefteq\pi\), write \(a\equiv b\pmod{\Gamma_{k+1}}\) when \(a\Gamma_{k+1}=b\Gamma_{k+1}\) in \(\pi/\Gamma_{k+1}\). With the intended groups and actions fixed, the requested assertion is
\[
\begin{gathered}
\forall\varphi\in A_{g,1}(k)\ \forall\psi\in B_{g,1},\\
\left[\forall x\in\pi,\quad
\varphi(x)\equiv\psi(x)\pmod{\Gamma_{k+1}}\right]\\
\Longrightarrow
\left[\exists\mu\in AB_{g,1}\ \forall x\in\pi,\quad
\mu(x)\equiv\varphi(x)\equiv\psi(x)\pmod{\Gamma_{k+1}}\right].
\end{gathered}
\]
This is to be considered for the intended admissible genera \(g\) and filtration indices \(k\); the record explicitly specifies \(k\geq1\) for the first question but does not separately specify the parameter ranges for the second.

\emph{Unresolved qualifications.} The record does not identify \(\pi\), define \(\Gamma_\bullet\) or its indexing, specify the relevant Johnson filtration and handlebody subgroups, or explain the notation \(AB_{g,1}\). In particular, its meaning as an intersection, product, or separately defined subgroup must not be presumed. The surface, boundary and basepoint conventions needed to interpret the actions on \(\pi\) are also absent. Consequently, these formulas preserve the supplied targets but do not constitute a fully specified mathematical formulation without the missing definitions.
\subsection{Short English statement}
Describe the Johnson-filtration images \(V_k(3)\) associated with stable Heegaard splittings of integral homology \(3\)-spheres, and ask whether every such image equals \(V_1(3)\). Also ask whether the common action modulo \(\Gamma_{k+1}\) of the specified elements \(\varphi\) and \(\psi\) can be realized by one element \(\mu\in AB_{g,1}\). Several essential group and filtration definitions are missing from the supplied record, so the formulation requires human mathematical review.
\subsection{Sources}
\begin{itemize}
\item Local UnsolvedMath registry, permanent ID \(30003289\), exact source problem code \texttt{catalogue entry}, title: \emph{Johnson-Filtration Images of Stable Heegaard Splittings}; collection: Oberwolfach Reports Open Problems. The supplied proposed year is 2016; no individual proposer is identified.
\item Supplied report citation: \emph{Surface Bundles (2017)}, Oberwolfach Reports, pp.\ 3149--3195, DOI \texttt{10.4171/owr/2016/56}. \url{https://doi.org/10.4171/owr/2016/56}. The supplied background attributes the problem to workshop open-problem discussions documented in Oberwolfach Reports; no problem-specific page or passage is supplied.
\item Supplied literature source: \emph{Oberwolfach Report 56 (2017)}. \url{https://doi.org/10.4171/owr/2016/56}. Literature checked at: \texttt{2026-08-21}. This is supplied dated curation, not an independently verified assessment. It labels the problem open and states that Johnson-filtration and mapping-class-group methods give partial information, while the asserted equality of the ``third-level images'' and the requested simultaneous realization remain unresolved. The supplied background explicitly calls for independent verification of time-sensitive status and cited claims.
\item The registry labels the statement \texttt{corrected\_verified} and reports manual semantic, typographic, and LaTeX review, including a claim of self-containment. These are supplied editorial assertions; the missing definitions identified here remain unresolved, and no source visit or independent mathematical verification is asserted.
\end{itemize}
\end{document}
